\documentclass[11pt]{article}
\usepackage{fontspec}
\usepackage[french]{babel}
\usepackage{amsmath}
\usepackage{amssymb}
\usepackage{amsthm}
\usepackage{mathrsfs}
\usepackage{tikz-cd}
\usepackage{enumitem}
\usetikzlibrary{babel}

\newcommand{\printedpage}[1]{\marginpar{\footnotesize p.~#1}}
\newcommand{\msp}[1]{\marginpar{\footnotesize [#1]}}
\newcommand{\margmark}[1]{\marginpar{\footnotesize #1}}
\newcommand{\PS}{\mathrm{PS}}
\newcommand{\Ob}{\mathrm{Ob}}
\newcommand{\tl}[1]{\begin{tabular}{@{}l@{}}#1\end{tabular}}
\newcommand{\tls}[1]{\text{\footnotesize\begin{tabular}{@{}l@{}}#1\end{tabular}}}
\newcommand{\circled}[1]{\tikz[baseline=(c.base)]\node[draw,circle,inner sep=1pt](c){#1};}

\begin{document}

%======================================================================
\printedpage{59}
% head: A. Grothendieck : Dérivateurs, Ch. II, Édition des Universités de Montpellier II et Paris VII

ouverte. Comme $y\backslash X$ est localement cofiltrant (resp. localement faiblement cofiltrant) par hypothèse, il en est donc de même de $y\backslash X$. Cela prouve d).

Appliquant a) à la situation $X/x \longrightarrow Y/y$ dans b), on voit que les colocalisés $z\backslash(X/x)$ de
\margmark{!} % margin: "!" with a triple vertical rule alongside the lines "X/x sur Y/y ... donc est 0-connexe."
$X/x$ sur $Y/y$ ($z \in \Ob\,(Y/y)$) sont localement cofiltrants (resp. faiblement). Mais comme $X/x$ a un objet final au dessus de l'objet final de $Y/y$, il s'ensuit que $z\backslash(X/x)$ a un objet final, donc est $0$-\textit{connexe.}

\msp{page 95}
Donc ces catégories, étant localement cofiltrantes (resp. faiblement), sont en fait cofiltrantes (resp. faiblement cofiltrantes).

Prouvons c). Notons que pour $z$ dans $V$, $z\backslash U$ s'identifie à l'image inverse, par $z\backslash(X/x) \longrightarrow X/x$, de l'ouvert $U$ de $X/x$, c'est donc un \textit{ouvert} de $z\backslash(X/x)$, non vide par construction de $V$. Mais un ouvert non vide d'une catégorie cofiltrante (resp. faiblement cofiltrante) est encore cofiltrant (resp. faiblement cofiltrant). Cela signifie que $g : U \longrightarrow V$ est cofiltrant (resp. faiblement cofiltrant), a fortiori cofinal (et $W$-coasphérique, pour tout $W$ satisfaisant Loc (8)). Cela achève la démonstration.

\noindent\textbf{Remarque 7.5.}
\begin{itemize}
\item[1)] Le même argument que dans c) montre que si $f : X \longrightarrow Y$ est cofiltrant (resp. faiblement cofiltrant), alors pour tout ouvert $U$ dans $X$, si on désigne par $V$ l'ouvert de $Y$ engendré par $f(U)$, i.e. l'ouvert formé des $y$ dans $Y$ tels que $y\backslash X \neq \emptyset$, alors la flèche induite $g : U \longrightarrow V$ est cofiltrante (resp. faiblement cofiltrante), et a fortiori $g$ est cofinal (et plus généralement, si $f$ est cofiltrant, $g$ est $W$-coasphérique pour tout $W$ satisfaisant la condition Loc (8)).

\msp{page 96}
\item[2)] Un cas particulier intéressant de 7.4 est celui où on suppose que les $y\backslash X$ ($y \in \Ob\,Y$) sont sommes directes de catégories ayant chacune un \textit{objet initial.} Alors on voit que cette condition est stable par les opérations de 7.4 a) [\,?], et par conséquent, les colocalisés $z\backslash(X/x)$ des morphismes localisés en haut $X/x \longrightarrow Y/y$ sont des catégories à objet initial. Ceci implique alors qu'il en est de même des colocalisés $z\backslash U$ ($z \in \Ob\,V$) dans $g : U \longrightarrow V$. Il s'ensuit que $g$ est $W$-coasphérique pour \textit{tout} localiseur fondamental $W$ (sans avoir à se soucier de l'axiome des limites).

La condition envisagée sur $f : X \longrightarrow Y$ est satisfaite notamment dans les deux cas suivants (qui me paraissent de loin les plus importants) : a) $f$ est étale et b) $X$ est discrète.
\end{itemize}
% foot: Version 8 juillet 2001

%======================================================================
\clearpage
\printedpage{60}
% head: A. Grothendieck : Dérivateurs, Ch. II, Édition des Universités de Montpellier II et Paris VII

\msp{page 97}

\begin{center}
{\small
\begin{tikzcd}[ampersand replacement=\&, column sep=small, row sep=large]
\left\{\begin{tabular}{@{}c@{\ }l@{}}
 & Top $X$\\
 & localement totalement 0-connexe\\
 & et 0-connexe\\
$\Longleftrightarrow$ & totalement 0-connexe\\
 & et localement irréductible\\
$\Longleftrightarrow$ & $X$ \textit{cofiltrante,}\\
 & i.e. $\PS^*1' + \PS^*2$\\
$\Longleftrightarrow$ & $X$ 0-connexe\\
 & $+ \PS^*1 + \PS^*2$\\
$\Longleftrightarrow$ & $\left\{\begin{tabular}{@{}l@{}}
$\forall\, x \in \Ob\,X$, $X/x$ et\\
$X/x \times x$ sont irréductible,\\ % [sic] irréductible -> irréductibles, agreement with "sont"
$X$ 0-connexe
\end{tabular}\right.$\\
$\Longleftrightarrow$ & $\left\{\begin{tabular}{@{}l@{}}
$X$ 0-connexe,\\
les $X/a \times_x b$ non vides\\
(et même 0-connexes),\\
les $X/a \times_{x\times x} b$ 0-connexes
\end{tabular}\right.$
\end{tabular}\right.
\;
\left[\begin{tabular}{@{}l@{}}
stable par\\
passage\\
à $X' \longrightarrow X$\\
étale sur $X$, si\\
$X'$ 0-connexe,\\
\textit{mais} si les\\
$X/x$ sont cofiltrantes\\
et $X$ 0-connexe,\\
cela n'implique\\
\textit{pas} $X$ cofiltrante\\
Ex. $B_G$
\end{tabular}\right]
\arrow[rrrd, Rightarrow, "{\tl{stricte\\(ex. $\bullet\ \bullet$)}}"]
\arrow[rrdd, Rightarrow, "{\tl{stricte\\(ex. $\boldsymbol{\Delta}$)}}"]
\arrow[ddd, Rightarrow, "{\tl{stricte (ex.\\$B_G$)}}"']
\& \& \& |[draw,rounded corners=6pt]| \varepsilon = \bullet\ \bullet \\
\& \& \&
\begin{tabular}{@{}l@{}}
\begin{tabular}{@{}c@{\ }l@{}}
 & Top $X$ localement\\
 & totalement 0-connexe\\
$\Longleftrightarrow$ & les composantes connexes de $X$\\
 & sont cofiltrantes\\
$\Longleftrightarrow$ & $\PS^*1 + \PS^*2$\\
 & (X \textit{localement cofiltrante}) % [sic] X (upright roman) -> $X$ (italic), inconsistent with all other occurrences
\end{tabular}\\[2ex]
$\left[\begin{tabular}{@{}l@{}}
Stable par passage à $X' \longrightarrow X$\\
étale, mais si les $X'/x$\\
y satisfont, pas forcément $X'$!
\end{tabular}\right]$
\end{tabular}
\arrow[d, Rightarrow] \\
\& \&
\begin{tabular}{@{}l@{}}
\begin{tabular}{@{}c@{\ }l@{}}
 & Top $X$ totalement 0-connexe\\
 & (on dit $X$ \textit{totalement} 0-\textit{connexe})\\
$\Longleftrightarrow$ & $X^o \longrightarrow X^o \times X^o$ est cofinal\\
 & \textit{et} $X \neq \emptyset$\\
$\Longleftrightarrow$ & $\forall\, x, y \in \Ob\,X$, $x \times y$ 0-connexe de $X^\wedge$\\
 & i.e. $X/x \times y$ 0-connexe\\
 & \textit{et} $X \neq \emptyset$
\end{tabular}\\[2ex]
$\left[\begin{tabular}{@{}l@{}}
NB n'est pas stable par passage aux $X/x$ (\footnotemark[66]), et si\\
les $X/x$ sont totalement 0-connexes et $X \neq \emptyset$,\\
il ne s'ensuit pas que\\
$X$ le soit (\footnotemark[67]).
\end{tabular}\right]$
\end{tabular}
\arrow[dddl, Rightarrow, "{\tl{stricte (ex. $B_G$, ou\\$0 \rightrightarrows 1$)}}"]
\&
\tl{Top $X$\\localement irréductible} \\
\begin{tabular}{@{}c@{\ }l@{}}
 & Top $X$ localement\\
 & irréductible\\
 & $+$ 0-connexe\\
$\Longleftrightarrow$ & $X$ \textit{pseudofiltrante}\\
 & $+$ 0-\textit{connexe,}\\
 & i.e. $\PS^*1$\\
 & $+$ 0-connexe
\end{tabular}
\arrow[ddr, Rightarrow, "{\tl{stricte\\(ex. $0 \rightrightarrows 1$)}}"']
\&
\left[\begin{tabular}{@{}l@{}}
$\PS^*1$ ne dépend\\
que des catégories\\
induites $X/x$,\\
donc stable par passage\\
aux $X' \longrightarrow X$ étales,\\
donc la condition\\
$\PS^*1 +{}$0-connexe stable\\
par passage aux\\
revêtements\\
$X' \longrightarrow X$ 0-\textit{connexes.}
\end{tabular}\right]
\& \& \\
|[draw,rounded corners=6pt]| B_G \& \& \& \\
\&
\begin{tabular}{@{}c@{\ }l@{}}
 & Top $X$ irréductible\\
 & ($X$ \textit{irréductible})\\
$\Longleftrightarrow$ & $\PS^*1'$
\end{tabular}
\& \& \\
|[draw,rounded corners=6pt]| 0 \rightrightarrows 1
\&
\left[\begin{tabular}{@{}l@{}}
ne dépend que \texttt{[de]} $\mathrm{Ord}\,X$ ou de $\mathrm{Esp}(X)$,\\
signifie que $\mathrm{Ord}\,X$ cofiltrant,\\
ou $\mathrm{Esp}(X)$ irréductible.\\
N'est pas stable par passage à\\
$X' \longrightarrow X$ étale 0-connexe,\\
p.ex. par passage à un $X/x$,\\
p.ex. $X = (0 \rightrightarrows 1)$ est irréductible,\\
mais $X/1 = \left(\begin{tabular}{@{}l@{}}$\cdot$\\ $\;\searrow$\\ $\qquad\cdot$\\ $\;\nearrow$\\ $\cdot$\end{tabular}\right)$ ne l'est\\
pas.
\end{tabular}\right]
\& \& \\
\end{tikzcd}
}
\end{center}
% footnote marks (66), (67) are set inside the diagram; texts below
\footnotetext[66]{\textbf{Ex.} \circled{$\boldsymbol{\Delta}$}.}
\footnotetext[67]{\textbf{Ex.} $B_G$.}
% foot: Version 8 juillet 2001

%======================================================================
\clearpage
\printedpage{61}
% head: A. Grothendieck : Dérivateurs, Ch. II, Édition des Universités de Montpellier II et Paris VII

\msp{page 98}

\begin{center}
{\small
\begin{tikzcd}[ampersand replacement=\&, column sep=small, row sep=large]
\& \tl{Top $X$ localement totalement 0-connexe\\$+$ 0-connexe\\($X$ \textit{cofiltrante})}
\arrow[ddl, Rightarrow] \arrow[ddr, Rightarrow] \arrow[rrd, Rightarrow]
\& \& \\
\& \& \& |[draw,rounded corners=6pt]| \varepsilon = \bullet\ \bullet \\
\tl{Top $X$ localement irréductible\\$+$ 0-connexe\\($X$ \textit{pseudofiltrante} et 0-\textit{connexe})\\$\PS^*1 +{}$ 0-connexe}
\arrow[ddr, Rightarrow]
\& \&
\tl{Top $X$ totalement 0-connexe\\($X$ \textit{totalement} 0-\textit{connexe})\\$X^o \longrightarrow X^o \times X^o$ cofinal}
\arrow[ddl, Rightarrow]
\&
\tl{Top $X$ localement totalement\\0-connexe\\($X$ \textit{localement cofiltrante})\\$\PS^*1 + \PS^*2$}
\arrow[dd, Rightarrow] \\
|[draw,rounded corners=6pt]| B_G \& \& |[draw,circle,inner sep=1pt]| \boldsymbol{\Delta} \& \\
\& \tl{Top $X$ irréductible\\($X$ \textit{irréductible})\\$\PS^*1'$}
\& \&
\tl{Top $X$ localement irréductible\\($X$ \textit{pseudo-cofiltrante} [\,?])\\$\PS^*1$} \\
\& |[draw,rounded corners=6pt]| 0 \rightrightarrows 1 \& \&
|[draw,rounded corners=12pt]| \begin{tabular}{@{}c@{}}$B_G \amalg e$\\avec $G \neq 1$\end{tabular} \\
\end{tikzcd}
}
\end{center}

\begin{center}
{\small
\begin{tikzcd}[ampersand replacement=\&, column sep=small, row sep=large]
\text{(\footnotemark[68]).}
\& \tl{$X$ \textit{cofiltrante}\\$\PS^*1' + \PS^*2$\\$= \PS^*1 + \PS^*2 +{}$ 0-connexe}
\arrow[ddl, Rightarrow] \arrow[ddr, Rightarrow] \arrow[rrd, Rightarrow]
\& \& \\
\& \& \& |[draw,rounded corners=6pt]| \bullet\ \bullet \\
\tl{$X$ \textit{pseudo-cofiltrante} et 0-\textit{connexe}\\($\PS^*1 +{}$ 0-connexe)}
\arrow[ddr, Rightarrow] \arrow[ddrrr, Rightarrow]
\& \&
\tl{$X$ \textit{totalement} 0-\textit{connexe}\\($X^o \longrightarrow X^o \times X^o$ cofinal)}
\arrow[ddl, Rightarrow]
\&
\tl{$X$ \textit{localement cofiltrante}\\$\PS^*1 + \PS^*2$}
\arrow[dd, Rightarrow] \\
|[draw,rounded corners=6pt]| B_G \& \& |[draw,circle,inner sep=1pt]| \boldsymbol{\Delta} \& \\
\& \tl{$X$ \textit{irréductible}\\$\PS^*1'$}
\& \&
\tl{$X$ \textit{pseudo-cofiltrante}\\$\PS^*1$} \\
\& |[draw,rounded corners=6pt]| \bullet \rightrightarrows \bullet \& \&
|[draw,rounded corners=6pt]| B_G \amalg e \\
\end{tikzcd}
}
\end{center}
\footnotetext[68]{\textbf{NB} Les cinq catégories typiques $B_G$ ($G \neq 1$), $\boldsymbol{\Delta}$, $0 \rightrightarrows 1$, $\varepsilon = \bullet\ \bullet$, $B_G \amalg e$ ($G \neq 1$) sont telles qu'elles appartiennent à la classe envisagée $C$ d'objets de Cat, mais à aucune des cinq autres classes sauf celles qui sont reliées à $C$ par $C \Longrightarrow C'$ sur le diagramme d'implications (ou par une implication composée).}
% foot: Version 8 juillet 2001

%======================================================================
\clearpage
\printedpage{62}
% head: A. Grothendieck : Dérivateurs, Ch. II, Édition des Universités de Montpellier II et Paris VII

\noindent\begin{tabular}{@{}p{2.8cm}@{}p{11.5cm}@{}}
$\bullet \rightrightarrows \bullet$ & n'est ni pseudo-cofiltrante, ni totalement 0-connexe.\\[1ex]
$B_G$ & n'est ni localement cofiltrante, ni totalement 0-connexe.\\[1ex]
$\boldsymbol{\Delta}$ & n'est pas pseudo-filtrante.\\[1ex]
$B_G \amalg e$ & n'est pas localement cofiltante, car $B_G$ ne l'est pas, ni appartient à une des 4 classes d'objets connexes.\\[1ex] % [sic] cofiltante -> cofiltrante, letter r missing
Itou pour & $\varepsilon = \bullet\ \bullet$
\end{tabular}

\msp{page 99}

\begin{center}
{\small
\begin{tikzcd}[ampersand replacement=\&, column sep=small, row sep=large]
\& \tl{$X$ \textit{cofiltrante}\\($\PS^*1' + \PS^*2$\\$\Longleftrightarrow \PS^*1 + \PS^*2$\\$+$ 0-connexe)}
\arrow[ddl, Rightarrow, "{\tls{$\PS^*1'$\\($\Longrightarrow$ 0-connexe)}}"]
\arrow[dd, Rightarrow, "{\tls{$\PS^*1$\\(les $x\backslash X$ satisfont\\$\PS^*1'$)}}"]
\arrow[dr, Rightarrow]
\& \\
% layout: the two small labels stand between the arrows; attribution to the arrows as in this layout
\& \& |[draw,rounded corners=6pt]| \varepsilon = \bullet\ \bullet \\
\tl{$X$ totalement 0-connexe\\($X^o \longrightarrow X^o \times X^o$ cofinal)}
\arrow[ddd, Rightarrow]
\&
\tl{$X$ \textit{pseudo-cofiltrante}\\0-connexe\\($\PS^*1 +{}$ 0-connexe)\\(manque $\PS^*2$)}
\arrow[dddl, Rightarrow] \arrow[dddr, Rightarrow]
\&
\tl{$X$ \textit{localement cofiltrante}\\(composantes connexes filtrantes)\\($\PS^*1 + \PS^*2$)}
\arrow[ddd, Rightarrow] \\
|[draw,circle,inner sep=1pt]| \boldsymbol{\Delta} \& |[draw,rounded corners=6pt]| B_G \& \\
\& G \neq \{1\} \& \\
\tl{$X$ irréductible\\($\PS^*1'$)}
\& \&
\tl{$X$ \textit{pseudo-cofiltrante}\\ou encore \textit{localement irréductible}\\(composantes connexes\\pseudo-filtrantes)\\($\PS^*1$)} \\
|[draw,rounded corners=6pt]| \bullet \rightrightarrows \bullet \& \& |[draw,rounded corners=6pt]| B_G \amalg e \\
\& \& G \neq \{1\} \\
\end{tikzcd}
}
\end{center}

et dualement on a

\begin{center}
{\small
\begin{tikzcd}[ampersand replacement=\&, column sep=small, row sep=huge]
\& \text{$X$ filtrante}
\arrow[dl, Rightarrow] \arrow[d, Rightarrow] \arrow[dr, Rightarrow]
\& \\
\text{$X^o$ totalement 0-connexe}
\arrow[d, Rightarrow]
\&
\tl{$X$ pseudo-filtrante\\0-connexe}
\arrow[dl, Rightarrow] \arrow[dr, Rightarrow]
\&
\text{$X$ localement filtrante}
\arrow[d, Rightarrow] \\
\text{$X$ coirréductible} \& \& \text{$X$ pseudo-cofiltrante} \\
\end{tikzcd}
}
\end{center}
% foot: Version 8 juillet 2001

%======================================================================
\clearpage
\printedpage{63}
% head: A. Grothendieck : Dérivateurs, Ch. II, Édition des Universités de Montpellier II et Paris VII

donc on a une terminologie directe sauf pour `$X^o$ totalement 0-connexe' - je n'ose dire `$X$ totalement co-connexe'!!
% foot: Version 8 juillet 2001

\end{document}
